\section{Final Discussion}
\label{sec:discussion}

\textbf{What the data taught us.} Three things shaped the whole project. First, the rare classes are not spread over patients: most F beats come from one training record and three quarters of the DS2 S beats from one test record (232), so ``learning S'' and ``learning F'' largely means learning one person, and the inter-patient score measures how well that transfers to another person, which is poorly. Second, the information that separates an S beat from an N beat is timing, not shape; the best single change in the project was widening the input window so that the network could see the neighbouring beats (validation S recall doubled), and even then timing failed on a patient whose rhythm does not follow the usual pattern. Third, the choice of validation records matters as much as the model: with the noisiest record and both bundle-branch-block records in validation, model selection chose an almost untrained network; after the revision of Section~\ref{sec:record-selection} the sweep became informative.

\textbf{The biggest source of error} is the S class, in both directions: 1\,550 S beats called N and 1\,488 N beats called S on DS2. One patient (232) supplies most of the first group; an atrial-fibrillation patient (219), whose normal beats arrive irregularly, half of the second. Together they cost more macro-F1 than everything else combined, and they explain the gap between validation (0.67) and DS2 (0.49). N and V, the classes with a shape signature, are handled well on unseen patients.

\textbf{Was the highest-accuracy model also the best Edge-AI model?} Not quite, and the difference is instructive. The random forest has the highest DS2 accuracy (0.935) but the lowest macro-F1, because it predicts N almost always and finds 0.3\,\% of the S beats. The plain cross-entropy CNN had the highest validation accuracy but a poor S recall. The final model trades about 1.5 points of accuracy for V recall of 0.94 and whatever S recall the data allowed, and it is the one that runs on the chip: 25\,kB of weights, five accelerator layers, no operator on the CPU. Among the CNNs, size did not matter: the 7\,428-parameter version was within 0.07 of the final on validation and would have been the right choice if memory had been tight.

\textbf{How much performance was lost after quantisation?} 0.012 macro-F1 on DS2, after quantisation-aware training with per-layer activation scales. Without that preparation the loss would have been about 0.27: the FP32 network uses activation ranges the 8-bit clamp cuts off. The lesson is that for this accelerator QAT is not an optimisation but a requirement, and that the activation-threshold step of the toolchain is the part that does the work.

\textbf{How sensitive was the system to noise?} Less than expected for baseline wander and muscle artefact, which the median filter and the network absorb (drop at 6\,dB below 0.02 on average), and clearly sensitive to electrode motion, which imitates ventricular beats (macro-F1 0.405 at 0\,dB, V F1 from 0.90 to 0.59). The low-pass filter turned out to be unnecessary. Robustness came from the preprocessing and the augmentation, not from model size; the largest network was the least robust.

\textbf{What compromises did the hardware force?} The operator set (no attention, no recurrence, no normalisation at inference) was not a real limit for this problem; a small 1D CNN is the right tool anyway. Two constraints did shape decisions: the fully connected layer may take at most 1\,024 inputs, which fixed the pooling schedule of the wide-window model, and the 8-bit activation range, which forced QAT and the activation scaling. The fixed window size means the model cannot use variable-length context such as the last ten RR intervals, which is exactly what would help the S class; adding those as a second small input is possible on the accelerator but was not attempted.

\textbf{What we would change with more time.} In order of expected gain: (1) give the network the RR features explicitly, as a second input branch merged before the output layer, so that timing does not have to be read from the waveform; (2) train with real \texttt{nstdb} noise at random SNR, which targets the electrode-motion weakness directly; (3) evaluate with record-wise cross-validation inside DS1 instead of one validation split, to make model selection less dependent on five records; (4) a per-patient adaptation step on the device, using the first minute of a recording to calibrate the N class, which published inter-patient work shows to be the largest single improvement for S; and (5) a real board, to turn the software-verified claims of Section~\ref{sec:deployment} into measurements.
